\begin{abstract}
Course-specific LLM assistants usually treat each student turn independently or remember only raw chat history, so they cannot tell \emph{which} prerequisite a student is missing or \emph{which} misconception keeps causing errors. We present LEARN, a course copilot whose central component is a structured, concept-level learner state: Bayesian Knowledge Tracing mastery probabilities and explicit misconception sets defined over an instructor-curated prerequisite graph. This state is updated after every turn by a diagnoser and directly controls four behaviors: information-gain-driven prerequisite probing, state-conditioned hybrid retrieval, mastery-scaled Socratic hints, and next-action selection. Responses are grounded in course PDFs and lecture transcripts through hybrid dense+BM25 retrieval with enforced citations and a grounding Critic.
We evaluate with (i) a component ablation (no RAG $\rightarrow$ vector RAG $\rightarrow$ hybrid cited RAG $\rightarrow$ +Critic $\rightarrow$ +learner state) on a \todo{N}-item single-turn benchmark and (ii) \todo{N} controlled multi-turn simulated-student trajectories with hidden prerequisite gaps and persistent misconceptions, comparing no memory, raw chat history, summary memory, and structured state.
\todo{Key results: e.g., structured state improves gap-identification F1 from X to Y and reduces unnecessary remediation by Z\% versus raw history; hybrid retrieval + Critic raises citation precision from A to B.}
We also report failure cases, including diagnoser errors that propagate through the state.
\todo{Keep to 150--250 words.}
\end{abstract}
